% Rayado de 9 mm (A4) — Eric Lucero González — https://ericlucero.dev
% Licencia CC BY-NC-SA 4.0. Compilar con pdfLaTeX (dos pasadas: latexmk lo hace solo).
\documentclass[a4paper]{article}
\usepackage[T1]{fontenc}
\usepackage[spanish]{babel}
\usepackage{lmodern}
\usepackage{tikz}
\pagestyle{empty}

% ── Configuración (mm) ─────────────────────────────────────
\def\anchoPapel{210}
\def\altoPapel{297}
\def\abajo{15}       % altura de la primera línea
\def\arriba{21}      % espacio libre mínimo arriba
\def\paso{9}        % distancia entre líneas
\def\margenIzq{30}   % línea de margen izquierda, desde el borde izquierdo
\def\margenDer{30}   % línea de margen derecha, desde el borde derecho

\begin{document}
\begin{tikzpicture}[
    remember picture, overlay, shift={(current page.south west)},
    x=1mm, y=1mm,
    normal lines/.style={gray!55, very thin},
    margin lines/.style={red!45!gray, thin}
  ]
  % Número de renglones que caben
  \pgfmathtruncatemacro{\n}{(\altoPapel - \abajo - \arriba) / \paso}

  \foreach \i in {0,...,\n} {
    \draw[normal lines] (0, \abajo + \i*\paso) -- (\anchoPapel, \abajo + \i*\paso);
  }

  % Líneas de margen
  \draw[margin lines] (\margenIzq, 0) -- (\margenIzq, \altoPapel);
  \draw[margin lines] (\anchoPapel - \margenDer, 0) -- (\anchoPapel - \margenDer, \altoPapel);
\end{tikzpicture}
\end{document}
